\chapter{Obesity and AI-Based Calorie Tracking}
\label{chap:obesity_calorie}

\section*{Introduction}
This chapter presents obesity as a major health issue, methods for managing it, principles of calorie tracking, and the role of artificial intelligence and computer vision in supporting dietary assessment........

\section{Obesity}

\subsection{Definition and Impact}
Overweight and obesity are defined based on Body Mass Index (BMI) and are associated with:
\begin{itemize}
    \item Health risks: cardiovascular diseases, diabetes
    \item Metabolic and hormonal disorders
    \item Reduced quality of life
\end{itemize}

\subsection{Management Approaches}
General strategies for managing obesity include:
\begin{itemize}
    \item Diet and calorie control
    \item Physical activity
    \item Medical supervision
\end{itemize}

\subsection{Calorie Tracking}
\begin{itemize}
    \item \textbf{Importance:} supports weight management and healthy habits
    \item \textbf{How it works:} monitors calorie intake versus energy expenditure
    \item \textbf{Tools:} mobile apps and websites for logging and tracking meals
\end{itemize}

\section{Artificial Intelligence for Calorie Tracking}

\subsection{Artificial Intelligence and Deep Learning}
\begin{itemize}
    \item Artificial Intelligence (AI): systems capable of performing intelligent tasks
    \item Deep Learning (DL): neural networks that learn patterns from data, especially images
\end{itemize}

Deep learning is widely used in several domains such as computer vision, healthcare, and food analysis. In the context of nutrition, it enables automated food recognition and calorie estimation from images.

\subsection{Food Estimation in Nutrition}
AI-based systems can analyze food images to support dietary assessment:
\begin{itemize}
    \item \textbf{Food recognition:} identifying food items in images
    \item \textbf{Calorie estimation:} analyzing meals to estimate their nutritional value
\end{itemize}

\subsection{Computer Vision Techniques}
Computer vision techniques are used to process and analyze food images:
\begin{itemize}
    \item \textbf{Image Classification:} CNN, MobileNet
    \item \textbf{Object Detection:} YOLO, Faster R-CNN
    \item \textbf{Segmentation:} Mask R-CNN
\end{itemize}

\section*{Conclusion}
This chapter introduced obesity, its impact and management, the role of calorie tracking, and the use of artificial intelligence and computer vision for food analysis and calorie estimation.....